% BHU PYQ -- History: History of India - Medieval (2024-25 NEP)
\documentclass[11pt,a4paper]{article}
\usepackage[a4paper,top=2.5cm,bottom=2.5cm,left=2.8cm,right=2.8cm,headheight=14pt]{geometry}
\usepackage{amsmath,amssymb}
\usepackage{fontspec}
\setmainfont{Kohinoor Devanagari}
\usepackage{microtype,fancyhdr,enumitem,hyperref,lastpage,parskip}

\hypersetup{colorlinks=true,urlcolor=black,linkcolor=black,pdftitle={HISMJ21: History of India - Medieval -- BHU NEP 2024-25}}

\pagestyle{fancy}\fancyhf{}
\renewcommand{\headrulewidth}{0.4pt}\renewcommand{\footrulewidth}{0.4pt}
\fancyhead[L]{\small \textit{Banaras Hindu University}}
\fancyhead[R]{\small \textit{HISMJ21: History of India - Medieval (2024-25)}}
\fancyfoot[C]{\small Page \thepage\ of \pageref{LastPage}}

\newlist{parts}{enumerate}{2}
\setlist[parts,1]{label=(\alph*),leftmargin=*,itemsep=4pt,topsep=3pt}
\setlist[parts,2]{label=(\roman*),leftmargin=*,itemsep=2pt,topsep=2pt}

\newcommand{\pts}[1]{\hfill \textit{[#1]}}

\begin{document}

\begin{center}
  {\large\bfseries BANARAS HINDU UNIVERSITY}\\[2pt]
  {\normalsize B. A. (Hons.) Semester II Examination, 2024-25 (NEP)}\\[6pt]
  \rule{0.8\linewidth}{0.5pt}\\[6pt]
  {\large\bfseries HISTORY}\\[4pt]
  {\large\bfseries Paper Code: HISMJ21 : History of India -- Medieval}\\[4pt]
  {\normalsize Time: 2:30 Hours \hfill Full Marks: 70}\\[2pt]
  {\small REF NO. CE/CONF./D-II/135}
\end{center}
\rule{\linewidth}{0.4pt}
\smallskip

\noindent \textit{Instructions: This question paper comprises three Sections A, B and C. All questions are compulsory. Write your Roll No.\ at the top immediately on receipt of this question paper.}

\smallskip
\noindent \textit{इस प्रश्न-पत्र में तीन खण्ड अ, ब तथा स हैं। सभी प्रश्न अनिवार्य हैं।}

\bigskip

\section*{SECTION -- A / खण्ड -- अ}
\noindent \textit{Note: Objective / Very Short Answer Questions. Answer each in not more than 50 words.}\pts{5 $\times$ 2 = 10}

\noindent \textit{वस्तुनिष्ठ / अति लघु उत्तरीय प्रश्न। प्रत्येक प्रश्न का उत्तर 50 शब्दों से अधिक नहीं होना चाहिए।}

\begin{enumerate}[label=\textbf{\arabic*.},leftmargin=*,itemsep=8pt]
  \item 
  \begin{parts}
    \item What is Paibos and Sijda?\pts{2}
    
    सिजदा एवं पाबोउस क्या हैं?

    \item Who was known as Tughril Tughan Khan / Tughril Beg?\pts{2}
    
    तुघ्राण खाँ / तुग्रिल बेग के रूप में किसे जाना जाता था?

    \item Who was the leader of the Turkish noble group during the Khalji Revolution?\pts{2}
    
    खिलजी क्रान्ति के समय तुर्क अमीर गुट का नेता कौन था?

    \item Who was Shahna-i-Mandi?\pts{2}
    
    शहना-ए-मंडी कौन था? (Market Superintendent under Alauddin Khalji)

    \item Write a short note on Diwan-i-Amir-i-Kohi.\pts{2}
    
    दीवान-ए-अमीर-ए-कोही पर एक टिप्पणी लिखिए। (Agriculture Department created by Muhammad bin Tughlaq)
  \end{parts}
\end{enumerate}

\bigskip

\section*{SECTION -- B / खण्ड -- ब}
\noindent \textit{Note: Short Answer Type Questions. Answer each question in about 250 words.}\pts{3 $\times$ 10 = 30}

\noindent \textit{लघु उत्तरीय प्रश्न। निम्नलिखित प्रत्येक प्रश्न का उत्तर लगभग 250 शब्दों में दीजिए।}

\begin{enumerate}[label=\textbf{\arabic*.},leftmargin=*,start=2,itemsep=12pt]

  \item Estimate the personality and character of Razia Sultana.\pts{10}
  
  रज़िया सुल्ताना के व्यक्तित्व एवं चरित्र का मूल्यांकन कीजिए।

  \begin{center}\textbf{OR / अथवा}\end{center}

  Discuss the Devgiri campaign of Alauddin Khalji.\pts{10}
  
  अलाउद्दीन खिलजी के देवगिरि अभियान की विवेचना कीजिए।

  \item Was the destruction of Hindu temples by Aurangzeb a part of his religious policy? Discuss.\pts{10}
  
  क्या औरंगज़ेब द्वारा हिन्दू मन्दिरों को तोड़ा जाना उसकी धार्मिक नीति का भाग था? विवेचना कीजिए।

  \begin{center}\textbf{OR / अथवा}\end{center}

  Describe the Mansabdari system introduced by Akbar.\pts{10}
  
  अकबर द्वारा प्रारम्भ की गई मनसबदारी व्यवस्था का वर्णन कीजिए।

  \item Discuss the religious policy of Iltutmish and his relation with the Caliphate.\pts{10}
  
  इल्तुतमिश की धार्मिक नीति एवं खलीफा के साथ उसके सम्बन्धों की विवेचना कीजिए।

  \begin{center}\textbf{OR / अथवा}\end{center}

  Write a note on the administrative reforms of Sher Shah Suri.\pts{10}
  
  शेरशाह सूरी के प्रशासनिक सुधारों पर एक टिप्पणी लिखिए।

\end{enumerate}

\bigskip

\section*{SECTION -- C / खण्ड -- स}
\noindent \textit{Note: Long Answer Type Questions. Answer the following questions in detail.}\pts{2 $\times$ 15 = 30}

\noindent \textit{दीर्घ उत्तरीय प्रश्न। निम्नलिखित प्रश्नों के विस्तृत उत्तर दीजिए।}

\begin{enumerate}[label=\textbf{\arabic*.},leftmargin=*,start=5,itemsep=14pt]

  \item Explain the Kingship theory of Ghiyasuddin Balban.\pts{15}
  
  बलबन के राजत्व सिद्धान्त की व्याख्या कीजिए।

  \begin{center}\textbf{OR / अथवा}\end{center}

  Describe the Rajput Policy of Akbar.\pts{15}
  
  अकबर की राजपूत नीति की व्याख्या कीजिए।

  \item Discuss the Economic Policy of Alauddin Khalji, including his market control policy.\pts{15}
  
  अलाउद्दीन खिलजी की आर्थिक नीति तथा बाज़ार नियंत्रण नीति की व्याख्या कीजिए।

  \begin{center}\textbf{OR / अथवा}\end{center}

  ``Muhammad bin Tughlaq was an idealist Sultan whose innovations harassed the people and ruined the empire.'' Critically examine this statement.\pts{15}
  
  ``मुहम्मद बिन तुग़लक़ एक विचारशील/आदर्शवादी सुल्तान था जिसकी नवीन योजनाओं ने जनता को परेशान किया और साम्राज्य को नष्ट कर दिया।'' इस कथन का समालोचनात्मक परीक्षण कीजिए।

\end{enumerate}

\end{document}
